\documentclass{ltjsbook}
\usepackage{docmute} 
\usepackage{../../myStyle} 

\begin{document}
\chapter{固有値と固有ベクトルと因子分析モデルの関係}
\label{x08_eigen}



\section{因子分析の数学的理解}
さあ因子分析に戻って考えてみましょう。因子分析のモデルは次のようなものでした。
\begin{equation}
    \bm{R} = \bm{A} \bm{A} ' + \bm{D}^2
\end{equation}

ここで，左辺の正方行列を相関行列$\bm{R}$とし，$\lambda_1 \bm{xx'} = \bm{aa'}$となるようにベクトルの大きさが整えられているとすれば，次のように表現できます。
\begin{equation}
    \bm{R}=\bm{a}_1\bm{a}_1'+\bm{a}_2\bm{a}_2'+\cdots +\bm{a}_m\bm{a}_m'+\bm{d} \bm{d}'
\end{equation}

ここでは共通因子の数が$m$個だとわかっている体で分解していますが，基本的にはサイズ$N$の行列からは$N$個の固有値がずらーっと並ぶわけです。それをどこかで「共通しているのはここまで」と判断し，残りは誤差であるとしてまとめて$\bm{dd'}$にしているだけですね。このように数学的にはここからが共通因子，ここからが独自因子といった区別をすることなく，最後のひとかけらまで固有値分解を行なっているのですが，その次元の重要度でもって共通因子と誤差因子に(研究者が恣意的に)分割しているのが因子分析のやっていることなのです。

一般に，$N$よりも$m$のほうがグッと少なくなります。たとえばYG性格検査では$N=120$であり，$m$はせいぜい5から10数個です。120項目のつくる120次元空間の中で，そこに働きかけても方向の変わらない基礎的な少数の次元にのみ注目すれば，効率よく情報圧縮ができるというものです。

相関係数を固有値分解すると，その固有値はすべて足し合わせるとサイズ$N$になるのでした。元のデータから計算される相関行列は，1つの項目が一単位分の情報を持っていると考えますが，固有値分解はそれを次元の重要性順に並べ替えます。固有値は項目いくつ分の重要度があるかということを表す指標だと考えることができます。どこから共通因子でどこからが誤差か，ということを考えるときに，たとえば固有値が1.0よりも小さくなるようであれば，項目1つ分の情報もないのだからということで誤差因子だと判断することがあります。

ところで因子分析モデルの$\bm{A}$，因子負荷行列ですが，これは共通因子の固有ベクトルをセットで扱ったものです。つまり，$\bm{A} = \bm{a_1,a_2,\cdots,a_m}$と縦ベクトルを並べたものになっています。エレメントワイズで表現すると次のようになります。

\[
    \bm{A} = \left (
    \begin{pmatrix}  a_{11}\\a_{12}\\\vdots a_{1N} \end{pmatrix},
    \begin{pmatrix}  a_{21}\\a_{22}\\\vdots a_{2N} \end{pmatrix},
    \cdots
    \begin{pmatrix}  a_{m1}\\a_{m2}\\\vdots a_{mN} \end{pmatrix},
    \right )
\]

ここで$\bm{AA'}$の間に単位行列$\bm{I}$ を挟んでも，別に結果は変わりませんよね。単位行列はかけても変わらないのが特徴ですから，$\bm{AA'}=\bm{AIA}$です。ここでの$\bm{I}$のサイズは$m \times m$であることに注意しつつ聞いてください。

$\bm{I} = \bm{TT'}=\bm{T'T}$になるような$m$次の行列を使うと，$\bm{AA'} = \bm{ATT'A'}$の関係は保たれたままです。この$\bm{T}$はこれまた行列の座標を変えてしまう変換行列であり，これを挟むことができるということは\textbf{因子負荷量の値はなんでもあり}だということになってしまいます。この変換行列$\bm{T}$はとくに\keyterm{回転行列}{かいてんぎょうれつ}{rotation matrix}とも呼ばれますが，因子分析はこのようにどんな回転でもできる，\textbf{回転に関する不定性}があるのです。あらゆる因子負荷量の組み合わせがあり得る，というのは困りものなので，なんらかの形で因子負荷行列$\bm{A}$に制約をかける必要があります。言い方を変えれば，色々な制約の中ではあっても好きな値を取ることができますから，ユーザにとって便利な基準を考えてやれば良いでしょう。因子分析では一般に，\keytermL{因子軸の回転}{いんしじくのかいてん}を行いますが，それはこうした理由からです。

ちなみに$\bm{TT'}$に
$diag(\bm{T\Phi T'})=\bm{I}_m$という
制約のある行列$\bm{\Phi}$を挟んでやっても，元の計算モデルに影響はありません。
この$\bm{\Phi}$は\keyterm{因子間相関}{いんしかんそうかん}{factor correlations}とよばれ，相関を持った因子軸の回転，すなわち\keyterm{斜交回転}{しゃこうかいてん}{oblique rotation}も色々なものが考えられています\footnote{これに対して$\bm{TT'}=\bm{I}$のような因子間相関がない(単位行列)な回転を\keyterm{直交回転}{ちょっこうかいてん}{orthogonal rotation}といいます。}。


ずいぶんややこしい話になってきました。この後は実践形式で因子分析を理解していければと思います。

\section{課題}
\paragraph{固有値問題}行列$\begin{pmatrix}8&1\\4&5\end{pmatrix}$の固有値・固有ベクトルを求めよ。


\end{document}
